\documentclass[12pt,a4paper]{article}
\usepackage{expediente}
\begin{document}
% =============================================================================
% FICHA DE DATOS DEL EXPEDIENTE --- actualizado 20 ago 2026
% =============================================================================

% --- VENDEDOR / TITULAR REGISTRAL -------------------------------------------
\newcommand{\VendedorNombres}{RICHAR DARWIN}
\newcommand{\VendedorApellidos}{CUTIRE CONDORI}
\newcommand{\VendedorNombreCompleto}{RICHAR DARWIN CUTIRE CONDORI}
\newcommand{\VendedorDNI}{42610690}
\newcommand{\VendedorEstadoCivil}{soltero, sin unión de hecho vigente}
\newcommand{\VendedorOcupacion}{\faltante{[ocupación de Richar]}}
\newcommand{\VendedorDomicilio}{\faltante{[domicilio actual de Richar]}}
\newcommand{\VendedorTelefono}{\faltante{[teléfono de Richar]}}
\newcommand{\VendedorCorreo}{\faltante{[correo]}}
\newcommand{\VendedorNacionalidad}{peruano}
\newcommand{\VendedorFechaNac}{\faltante{[fecha de nacimiento]}}
\newcommand{\VendedorDNIEmision}{\faltante{[fecha de emisión del DNI de Richar --- la pide SUNARP]}}

% --- CONVIVIENTE: NO CORRESPONDE -------------------------------------------
\newcommand{\ConvivienteNombre}{NO CORRESPONDE}
\newcommand{\ConvivienteDNI}{NO CORRESPONDE}
\newcommand{\ConvivienteDesde}{NO CORRESPONDE}
\newcommand{\ConvivienteInscrita}{no existe unión de hecho vigente}
\newcommand{\ConvivienteDomicilio}{NO CORRESPONDE}

% --- COMPRADORES (copropiedad en partes iguales) ---------------------------
\newcommand{\CompradorUnoNombre}{NICO ALVARO CUTIRE ARCE}
\newcommand{\CompradorUnoDNI}{48325017}
\newcommand{\CompradorUnoEmision}{01 de enero de 2024}
\newcommand{\CompradorUnoEstadoCivil}{soltero}
\newcommand{\CompradorUnoOcupacion}{ingeniero informático}
\newcommand{\CompradorUnoDomicilio}{Jr.\ Wiracocha, Sicuani, provincia de Canchis, Cusco, Perú}
\newcommand{\CompradorUnoTelefono}{984715108}

\newcommand{\CompradorDosNombre}{LUCY CUTIRE ARCE}
\newcommand{\CompradorDosDNI}{47478852}
\newcommand{\CompradorDosEmision}{09 de enero de 2026}
\newcommand{\CompradorDosEstadoCivil}{soltera}
\newcommand{\CompradorDosOcupacion}{contador}
\newcommand{\CompradorDosDomicilio}{Comunidad Hanocca, distrito de Layo, provincia de Canas, Cusco, Perú}
\newcommand{\CompradorDosTelefono}{953440875}

\newcommand{\CompradorNombres}{NICO ALVARO y LUCY}
\newcommand{\CompradorApellidos}{CUTIRE ARCE}
\newcommand{\CompradorNombreCompleto}{NICO ALVARO CUTIRE ARCE y LUCY CUTIRE ARCE}
\newcommand{\CompradorDNI}{48325017 y 47478852}
\newcommand{\CompradorEstadoCivil}{solteros}
\newcommand{\CompradorConyuge}{NO CORRESPONDE}
\newcommand{\CompradorConyugeDNI}{NO CORRESPONDE}
\newcommand{\CompradorOcupacion}{ingeniero informático y contador, respectivamente}
\newcommand{\CompradorDomicilio}{Sicuani (Canchis) y Layo (Canas), Cusco}
\newcommand{\CompradorTelefono}{984715108 / 953440875}
\newcommand{\CompradorCorreo}{\faltante{[correo de contacto]}}
\newcommand{\CompradorNacionalidad}{peruanos}
\newcommand{\PorcentajeCopropiedad}{50\,\% cada uno}

% --- INTERMEDIARIO / OCUPANTE DEL REMANENTE --------------------------------
\newcommand{\IntermediarioNombre}{LUCIO GIL CUTIRE PARI}
\newcommand{\IntermediarioDNI}{42467850}
\newcommand{\IntermediarioEstadoCivil}{\faltante{[estado civil de Lucio]}}
\newcommand{\IntermediarioOcupacion}{\faltante{[ocupación]}}
\newcommand{\IntermediarioDomicilio}{\faltante{[domicilio de Lucio]}}
\newcommand{\IntermediarioTelefono}{\faltante{[teléfono de Lucio]}}
\newcommand{\IntermediarioParentesco}{familiar (apellido Cutire)}

% --- PREDIO MATRIZ ----------------------------------------------------------
\newcommand{\PartidaMatriz}{\faltante{[N.º de partida electrónica SUNARP --- pendiente]}}
\newcommand{\OficinaRegistral}{Oficina Registral de Puerto Maldonado}
\newcommand{\AreaMatriz}{\faltante{[área total inscrita; las partes refieren del orden de 100 ha]}}
\newcommand{\UnidadCatastral}{\faltante{[código catastral / UC, si existe]}}
\newcommand{\NombreFundo}{\faltante{[nombre del fundo, si tiene]}}
\newcommand{\KmCarretera}{Interoceánica Sur, al este del Puente Jayave (el puente queda al oeste del predio, sobre el río; no debajo del lote)}

% --- PREDIO A INDEPENDIZAR --------------------------------------------------
\newcommand{\AreaIndependizarHa}{10{,}00}
\newcommand{\AreaIndependizarMetros}{100 000{,}00}
\newcommand{\PerimetroM}{2 200{,}58}
\newcommand{\FrenteM}{100{,}29}
\newcommand{\FondoM}{1 000{,}00}

\newcommand{\PUnoLat}{$-12{,}912452$}
\newcommand{\PUnoLon}{$-70{,}170058$}
\newcommand{\PDosLat}{$-12{,}912497$}
\newcommand{\PDosLon}{$-70{,}170981$}
\newcommand{\PTresLat}{$-12{,}903469$}
\newcommand{\PTresLon}{$-70{,}171438$}
\newcommand{\PCuatroLat}{$-12{,}903424$}
\newcommand{\PCuatroLon}{$-70{,}170515$}

\newcommand{\PUnoE}{373 065{,}00}
\newcommand{\PUnoN}{8 572 256{,}00}
\newcommand{\PDosE}{372 965{,}00}
\newcommand{\PDosN}{8 572 251{,}00}
\newcommand{\PTresE}{372 911{,}00}
\newcommand{\PTresN}{8 573 249{,}00}
\newcommand{\PCuatroE}{373 011{,}00}
\newcommand{\PCuatroN}{8 573 255{,}00}

\newcommand{\FechaPago}{15 de agosto de 2018}
\newcommand{\MontoPago}{S/\ 20 000{,}00}
\newcommand{\MontoPagoLetras}{VEINTE MIL Y 00/100 SOLES}
\newcommand{\FormaPagoHistorica}{dinero en efectivo, sin medio de pago del sistema financiero}

% Nombres genéricos --- sustituir por los reales al ir a notaría
\newcommand{\NotarioNombre}{Dr.\ JUAN CARLOS RÍOS SALINAS (nombre genérico a sustituir)}
\newcommand{\NotariaCiudad}{Puerto Maldonado}
\newcommand{\ColegiaturaAbogado}{CAL Madre de Dios N.º 0000 (genérico)}
\newcommand{\AbogadoNombre}{Abog.\ ANA MARÍA QUISPE HUAMÁN (nombre genérico a sustituir)}
\newcommand{\IngenieroNombre}{Ing.\ PEDRO ANTONIO MAMANI CONDORI (nombre genérico a sustituir)}
\newcommand{\IngenieroCIP}{CIP 000000 (genérico)}
\newcommand{\Municipalidad}{Municipalidad Distrital de Inambari}
\newcommand{\Provincia}{Tambopata}
\newcommand{\Departamento}{Madre de Dios}
\newcommand{\Distrito}{Inambari}

\newcommand{\LugarOtorgamiento}{Puerto Maldonado}
\newcommand{\DiaOtorgamiento}{\faltante{[día]}}
\newcommand{\MesOtorgamiento}{\faltante{[mes]}}
\newcommand{\AnioOtorgamiento}{2026}

\newcommand{\TestigoUno}{\faltante{[testigo 1: nombres, DNI, domicilio]}}
\newcommand{\TestigoDos}{\faltante{[testigo 2: nombres, DNI, domicilio]}}

\InicioDocumento{Declaración jurada de estado civil, convivencia y facultad de disponer}{DOC-12}

\noindent \textbf{SEÑOR NOTARIO PÚBLICO DE \MakeUppercase{\NotariaCiudad}}

\vspace{0.2em}
\noindent Yo, \textbf{\VendedorNombreCompleto}, de nacionalidad \VendedorNacionalidad, identificado con DNI N.º\ \textbf{\VendedorDNI}, de estado civil \textbf{\VendedorEstadoCivil}, de ocupación \VendedorOcupacion, con domicilio en \VendedorDomicilio, \textbf{declaro bajo juramento} lo siguiente:

\Clausula{Identidad y capacidad}
Soy mayor de edad, me encuentro en pleno ejercicio de mis derechos civiles y no tengo interdicción, ni medida judicial que me impida disponer de mis bienes. Mi DNI se encuentra vigente. Fecha de nacimiento: \VendedorFechaNac.

\Clausula{Estado civil}
Soy \textbf{soltero}. No contraje matrimonio civil ni religioso con efectos civiles que se encuentre vigente. Si el Registro de Personas Naturales o el DNI consignaren dato distinto, me obligo a subsanarlo antes de la escritura.

\Clausula{Unión de hecho / conviviente}
Declaro que \textbf{NO} mantengo unión de hecho. No convivo con pareja. Tuve hijos de una relación anterior; no nos casamos y ya no vivimos juntos. No existe conviviente que deba firmar. El DOC-17 recoge esta misma declaración.

\Clausula{Calidad propia o social del predio}
El predio rústico de la partida \PartidaMatriz, del cual se independizarán 10,00 ha, fue adquirido por mí:

\vspace{0.2em}
\noindent $\square$ \textbf{antes} de iniciar cualquier convivencia \quad $\square$ \textbf{durante} la convivencia \quad $\square$ por sucesión / otro: \linea[4.2cm]

\vspace{0.25em}
\noindent En consecuencia, para este acto de disposición:

\vspace{0.15em}
\noindent $\square$ considero el bien \textbf{propio} y no requiere intervención de conviviente, porque no hay unión de hecho vigente (DOC-17).

\Clausula{Descendientes}
Tengo hijos. Dejo constancia de que el acto proyectado es \textbf{compraventa a título oneroso} y no donación ni anticipo de legítima. Los descendientes \textbf{no intervienen} en la enajenación onerosa de bienes de su padre vivo. Me obligo a no simular donación.

\Clausula{Facultad de disponer}
Soy el único titular registral del predio matriz, según partida \PartidaMatriz. No he transferido, hipotecado ni prometido en exclusiva las 10 ha a persona distinta del comprador \CompradorNombreCompleto, salvo el trato informal de 2018 que se reconoce y se agota en este expediente. No existe copropietario inscrito. Si existiera cónyuge o conviviente con derecho, su consentimiento se otorga en documento aparte.

\Clausula{Veracidad}
Sé que la falsedad de esta declaración puede acarrear nulidad del acto, responsabilidad civil y las consecuencias penales de los artículos 411 y 427 del Código Penal, en cuanto correspondan.

\Clausula{Anexos}
Copia de DNI. $\square$ Ficha RENIEC. $\square$ Certificado negativo de unión de hecho. $\square$ Copia de inscripción de unión de hecho.

\LugarFecha

\noindent
\BloqueFirma{EL DECLARANTE / TITULAR}{\VendedorNombreCompleto}{\VendedorDNI}
\hfill
\BloqueFirmaSimple{CONVIVIENTE (enterada)}{\ConvivienteNombre\\DNI \ConvivienteDNI}

\vspace{0.8em}
\noindent{\small Certificación notarial de firmas (si se usa como documento separado):}\\
{\small Fecha: \linea[3cm] \quad Notario: \NotarioNombre\ \quad N.º de instrumento: \linea[3.2cm]}\\[0.3em]
\noindent\fbox{\parbox[c][2.2cm][c]{0.96\textwidth}{\centering\color{black!40}\small sello y firma del notario}}

\end{document}
